\begin{frame}
    \frametitle{\insertsubsectionhead：解码方法}
    \begin{block}{沿编码树逐位解码}
        \begin{enumerate}
            \item 发送方和接收方使用同一棵编码树，指针从根开始
            \item 依次读入一位：$0$走左子树，$1$走右子树
            \item 到达叶结点时输出对应符号，指针回到根，继续读取
            \item 完整码流结束时应回到根；若路径不存在或停在内部结点，需报错
        \end{enumerate}
    \end{block}
    \begin{exampleblock}{为何不用分隔符？}
        符号只位于叶结点，因此任何编码都不是另一个编码的前缀；
        到达叶结点即可确定一个符号。
    \end{exampleblock}
    \small 若实际字节末尾有填充位，须另行记录有效位数或符号数。
\end{frame}

\begin{fragile}
    \frametitle{\insertsubsectionhead：逐位解码示例}
    \bicolumns[0.45]{
        \begin{block}{码流：\texttt{\alert<2>{0}\alert<3>{0}\alert<4>{1}\alert<5>{1}%
                    \alert<6>{0}\alert<7>{1}\alert<8>{1}\alert<9>{0}}}
            \only<1>{从根开始，尚未读位。}
            \only<2>{读$0$，到结点$38$。}
            \only<3>{读$0$，到叶$A$；输出$A$，回到根。}
            \only<4>{读$1$，到结点$62$。}
            \only<5>{读$1$，到叶$B$；输出$B$，回到根。}
            \only<6>{读$0$，到结点$38$。}
            \only<7>{读$1$，到结点$20$。}
            \only<8>{读$1$，到结点$10$。}
            \only<9>{读$0$，到叶$F$；输出$F$，回到根。}
        \end{block}
        \begin{exampleblock}{已输出}
            \only<1-2>{空}
            \only<3-4>{$A$}
            \only<5-8>{$A,B$}
            \only<9>{$A,B,F$}
        \end{exampleblock}
    }{
        \begin{figure}
            \centering
            \begin{tikzpicture}[level distance=4ex, level/.style={sibling distance=23ex/#1}]
                \node[nonterminal] {\alert<1>{$100$}}
                    child[dotted] {node[nonterminal,solid] {\alert<2,6>{$38$}}
                        child[dotted] {node[terminal,solid] {\alert<3>{$A$}}}
                        child[solid] {node[nonterminal] {\alert<7>{$20$}}
                            child[dotted] {node[terminal,solid] {$C$}}
                            child[solid] {node[nonterminal] {\alert<8>{$10$}}
                                child[dotted] {node[terminal,solid] {\alert<9>{$F$}}}
                                child[solid] {node[terminal] {$G$}}
                            }
                        }
                    }
                    child[solid] {node[nonterminal] {\alert<4>{$62$}}
                        child[dotted] {node[nonterminal,solid] {$30$}
                            child[dotted] {node[terminal,solid] {$D$}}
                            child[solid] {node[terminal] {$E$}}
                        }
                        child[solid] {node[terminal] {\alert<5>{$B$}}}
                    };
            \end{tikzpicture}
            \caption{原示例编码树：虚线为$0$，实线为$1$}
            \label{fig:demo_huffman_decode}
        \end{figure}
    }
\end{fragile}

\begin{frame}
    \frametitle{\insertsubsectionhead：解码结果与边界}
    \begin{exampleblock}{与编码引例对应}
        $\texttt{00110110}=\texttt{00}\mid\texttt{11}\mid\texttt{0110}$
        解码为$A,B,F$；各段分别走到叶结点后才确定分界。
    \end{exampleblock}
    \begin{block}{实际使用时}
        \begin{itemize}
            \item 收发双方须共享或传输一致的码树或码表
            \item 未到叶结点就结束的有效码流是不完整的；不能把填充位当作有效数据
            \item 若字母表只有一个符号，需约定非空码字或另存符号数量
        \end{itemize}
    \end{block}
\end{frame}
